\chapter{Virtual Machines and Bytecode}\label{ch:virtual-machines-bytecode}

A virtual machine turns interpretation into an engineering discipline.
It fixes an instruction set, an encoding, and a state model before anyone writes an optimizer.
Programs become data that tools can verify, transform, and execute on any host.
This chapter defines a small bytecode machine precisely enough to implement and test.
We then show how verification, dispatch, and calls interact.
The result is a stable target for compilers and a stable host for later just-in-time techniques.

\section{Why Fix a Bytecode}\label{sec:ch31-why-bytecode}

A \term{virtual machine} is a program that executes programs written for an abstract instruction set.
The abstract set is called \term{bytecode} when instructions are compactly encoded.
Fixing the set separates language front ends from execution engines.
A front end lowers source constructs to a bounded list of operations.
The machine implements each operation once.
New languages reuse the machine, and new hosts reimplement it.
That reuse only works when the machine definition is complete.
It must specify value representations, instruction meanings, error behavior, and limits.
Consider unsigned addition on eight-bit values.
The definition must say whether the result wraps, traps, or produces a wider value.
It must say what happens on stack underflow and on a jump past the program end.
Without those answers, two hosts can run most programs yet disagree on edge cases.
We therefore define the machine as a state transition system.
A state holds a program counter, an operand stack, a frame stack, and a heap reference.
Each instruction maps one state to the next state or to a defined error.
The definition fits on a few pages for a small set.
Its brevity is a feature.
Every omitted case is a portability bug waiting for a test to find it.

\section{Stack Versus Register Machines}\label{sec:ch31-stack-register}

Bytecode machines divide into \term{stack machines} and \term{register machines}.
A stack machine keeps temporaries on an operand stack.
Instructions are short because operands are implicit.
A register machine names temporaries explicitly in each instruction.
Instructions are longer but programs need fewer of them.
Figure~\ref{fig:ch31-machine-state} shows the state components shared by both designs.
Control, operands, activation records, and heap references stay in distinct regions.
The split lets verification reason about each region separately.
Table~\ref{tab:ch31-stack-register} compares the two styles on code size, dispatch count, and verification effort.
The comparison is qualitative.
Real designs mix ideas, for example by adding a small register file to a stack machine.
The expression \code{(a + b) + (c + d)} illustrates the trade-off.
A stack encoding pushes four values and performs three additions, using seven instructions before return or halt.
If the four inputs are already in virtual registers, a register encoding can express the same calculation with three adds.
If separate loads are needed, those instructions must also be counted; the benefit depends on the chosen instruction set.
Neither style is uniformly better.
Stack code suits compact distribution and simple verification.
Register code suits fast dispatch and direct mapping to native registers.
We choose a stack machine for this chapter because its verifier is shorter.
Later chapters reuse the same state picture for frames and deoptimization.

\begin{figure}[htbp]
\centering
\begin{tikzpicture}[node distance=1.0cm and 1.2cm]
\node[basic block] (pc) {pc\\ next opcode};
\node[basic block, right=of pc] (op) {operand stack\\ temporaries};
\node[basic block, right=of op] (fr) {frames\\ locals, return pc};
\node[basic block, right=of fr] (hp) {heap\\ objects, arrays};
\draw[cfg edge] (pc) -- (op);
\draw[cfg edge] (op) -- (fr);
\draw[cfg edge] (fr) -- (hp);
\end{tikzpicture}
\caption{Machine state separates control, temporaries, activation records, and heap data so each instruction touches a defined region.}
\label{fig:ch31-machine-state}
\end{figure}

\begin{table}[htbp]
\centering
\caption{Stack and register designs differ in encoding density and dispatch count rather than expressiveness.}
\label{tab:ch31-stack-register}
\begin{tabular}{llll}
\toprule
Property & Stack machine & Register machine & Consequence \\
\midrule
Operands & Implicit & Explicit & Stack code is shorter \\
Dispatches & More & Fewer & Registers save steps \\
Verification & Simpler & Richer & Stack types track depth \\
Mapping & Indirect & Direct & Registers map to hardware \\
\bottomrule
\end{tabular}
\end{table}

\section{Encoding and Verification}\label{sec:ch31-encoding}

An encoding assigns each opcode a numeric tag and a fixed operand shape.
Fixed shapes keep decoding simple and verification decidable.
Our machine uses one-byte opcodes with zero, one, or two inline operands.
Jump targets are instruction indices, not byte offsets.
That choice removes alignment questions and makes out-of-range jumps easy to detect.
\term{Bytecode verification} checks a program before execution.
It proves that the operand stack has a consistent depth at every program point.
It proves that each instruction receives values of a compatible kind.
It proves that every jump target is a valid instruction start.
Algorithm~\ref{alg:ch31-verify} gives the depth-checking pass for straight-line and branching code without subroutines.
It assumes a control-flow graph over instruction indices is already built.
It requires each opcode to declare how many operands it pops and how many results it pushes.
It propagates a required depth along each edge and reports the first inconsistency.
Type checking extends the same idea from depths to kinds.
Each stack slot carries an abstract kind such as integer, reference, or unknown.
Handlers declare the kinds they consume and produce.
The verifier joins kinds at merge points and rejects programs with no consistent assignment.
Verification runs once per program load.
Execution then omits most static checks except guards the definition requires, such as division and null checks.
A frequent mistake verifies depth but skips kinds.
Depth-correct programs can still confuse integers and references.
Another mistake allows jumps into the middle of multi-word instructions.
Index-based targets prevent that error class by construction.
We adopt both restrictions here.

Figure~\ref{fig:visual-stack-arity} shows why the pop count cannot be replaced with the net stack effect.
An add started with one operand has net depth $0$, yet its second input never existed.
Listing~\ref{lst:visual-depth-step} checks the required operands and the remaining capacity in an order that avoids unsigned underflow and overflow in the verifier itself.

\begin{figure}[htbp]
  \centering
  \begin{tikzpicture}[node distance=13mm and 13mm]
    \node[basic block] (bad) {before: [7]};
    \node[basic block,right=of bad] (check) {ADD needs 2};
    \node[basic block,right=of check] (reject) {reject: only 1};
    \draw[cfg edge] (bad) -- (check);
    \draw[cfg edge] (check) -- (reject);
    \node[basic block,below=of bad] (good) {before: [7, 5]};
    \node[basic block,right=of good] (pop) {pop 2; push 1};
    \node[basic block,right=of pop] (accept) {after: [12]};
    \draw[cfg edge] (good) -- (pop);
    \draw[cfg edge] (pop) -- (accept);
    \node[font=\small,below=of pop] {$1-2+1=0$ is not proof that two operands existed.};
  \end{tikzpicture}
  \caption{An add with only one operand fails even though its net depth would be zero; the verifier must check pops before pushes.}
  \label{fig:visual-stack-arity}
\end{figure}



\begin{listing}[htbp]
  \inputminted{c}{example-code/visual-labs/depth.c}
  \caption{A checked depth transition validates operand availability before calculating the resulting depth.}
  \label{lst:visual-depth-step}
\end{listing}

The helper verifies one transition, not the whole control-flow graph or operand types.
The graph solver must still require matching depths at joins, valid targets, and the specified depth at each return.
Its tests enumerate 9,000 combinations of small depths, pop counts, push counts, and capacities, and separately exercise the largest \code{size\_t} boundary.

\begin{algorithm}[htbp]
\caption{Stack-depth verification over a bytecode control-flow graph}\label{alg:ch31-verify}
\begin{pseudocode}
\Require decoded array $P$, entry depth $d_0$, pop count $p(i)$ and push count $q(i)$
\Ensure depth map $D$ or a verification error
\State build CFG with successors for branches, jumps, and fall-through
\State set $D[v] \gets$ unknown for all $v$, $D[entry] \gets d_0$
\State $W \gets \{entry\}$
\While{$W \neq \emptyset$}
\State remove some $v$ from $W$
\If{$D[v] < p(v)$}
\State \Return underflow at $v$
\EndIf
\State $d \gets D[v]-p(v)+q(v)$
\State reject arithmetic overflow or $d$ above the configured stack capacity
\State reject invalid branch targets and invalid return depths
\ForAll{$s \in \Succ(v)$}
\If{$D[s]$ is unknown}
\State $D[s] \gets d$
\State add $s$ to $W$
\ElsIf{$D[s] \neq d$}
\State \Return depth mismatch at $s$
\EndIf
\EndFor
\EndWhile
\State \Return $D$
\end{pseudocode}
\end{algorithm}

Decode instruction boundaries and validate encoded targets before constructing successor links, including targets in unreachable code if the format requires whole-file validity.
The worklist verifies stack depths only on entry-reachable instructions; unreachable instructions retain unknown depth.
A successful depth check does not establish operand types or execution termination, which require separate contracts.

\section{Worked Example: Verifying and Running a Sum}\label{sec:ch31-worked}

We verify and execute a program that computes \code{(6 + 9) + (2 + 3)}.
The bytecode uses pushes, adds, and a halt.
Control flow is straight-line, so the verifier has one path to check.
Each push has effect $+1$ and each add has effect $-1$.
Starting from depth $0$, the depths evolve deterministically.
After the first push the depth is $1$.
After the second push the depth is $2$.
The first add reduces it to $1$.
The third push raises it to $2$.
The fourth push raises it to $3$.
The second add reduces it to $2$.
The final add reduces it to $1$.
Each add has at least two operands before it executes, and the exit depth is $1$.
Checking only that the net depth stays nonnegative would miss an add started with one operand.
Verification accepts with one value left for the halt.
The listing below records the expected depth and value at each step.
\begin{irlisting}
; (6 + 9) + (2 + 3), stack machine encoding
push 6        ; depth 0 -> 1
push 9        ; depth 1 -> 2
add           ; depth 2 -> 1, value 15
push 2        ; depth 1 -> 2
push 3        ; depth 2 -> 3
add           ; depth 3 -> 2, value 5
add           ; depth 2 -> 1, value 20
halt          ; result 20
\end{irlisting}
Execution proceeds from an empty stack.
After pushing $6$ and $9$, the stack holds $[6, 9]$.
The first add replaces them with $[15]$.
Pushing $2$ and $3$ gives $[15, 2, 3]$.
The second add gives $[15, 5]$.
The final add gives $[20]$.
The halt returns $20$.
Arithmetic confirms the result: $6 + 9 = 15$, $2 + 3 = 5$, and $15 + 5 = 20$.
There are seven value-producing instructions and one halt, for eight dispatches when halt is dispatched.
A register encoding would use fewer dispatches at the cost of wider instructions.
That is the density trade-off of Table~\ref{tab:ch31-stack-register} made concrete.
We also check two failure modes by hand.
Removing the final add leaves exit depth $2$, which a strict loader rejects when exactly one result is required.
Changing a jump target past the program end fails control-flow construction before depth checking begins.
Both checks happen before any instruction executes.

\section{Calls, Frames, and Constants}\label{sec:ch31-calls}

Calls need an explicit frame discipline on top of the operand stack.
A frame holds locals, a return address, and a saved stack base.
A call pushes a frame and resets the operand window.
A return restores both and delivers one result value.
Constants live in a separate pool so instructions stay narrow.
A load takes a pool index instead of embedding a full value.
That split keeps the common cases to one operand byte.
\begin{ccwarning}[One table]
The verifier, the interpreter, and any compiler must share one opcode table.
Separate tables drift apart and admit programs that one component rejects.
\end{ccwarning}
Test calls with three tiny programs.
The first calls a leaf and returns a constant.
The second nests two calls and checks operand windows do not leak.
The third passes the wrong arity and must fail verification, not execution.
The \enquote{fail before running} rule keeps calling-convention defects reproducible.

\section{Implementation Notes and Complexity}\label{sec:ch31-implementation}

A compact value stack backed by an array handles typical programs well.
Frames form a linked chain with a small header per activation.
The verifier of Algorithm~\ref{alg:ch31-verify} runs in $O(N + E)$ where $N$ is the instruction count and $E$ is the edge count.
Each program point enters the worklist at most twice: once when first assigned and once when a mismatch is found.
Memory use is $O(N)$ for the depth map and successor lists.
Dispatch cost per instruction is constant for a switch loop.
Total execution cost is $O(I)$ dispatches where $I$ is the dynamic instruction count, plus handler work.
Production machines add measured refinements.
They cache name lookups at call sites.
They tier from a switch loop to threaded code for hot methods.
They split verification into load-time structural checks and run-time type guards.
Each refinement preserves the state model of Figure~\ref{fig:ch31-machine-state}.
A useful discipline tests the verifier with mutants that shift one jump target or drop one push.
Those single-change mutants expose missing checks quickly.
Another discipline replays every verifier rejection as an interpreter test that must fail before dispatch.
% TODO(cite): need a primary reference for split bytecode verification designs.

\section{Keep adaptive handlers within the bytecode contract}
\label{sec:ch31-adaptive-link}

A verified bytecode operation can have several execution handlers without acquiring several language meanings.
An integer-specialized addition handler must fall back or report the specified error when its type guard fails.
A fused load/load/add handler must preserve the stack effect and attribute failures to the correct guest operation.
Verification of the original bytecode does not automatically validate a newly installed handler's assumptions.

Chapter~\ref{ch:adaptive-staged} implements a quickening state machine and a small syntax-specialization pipeline.
These examples keep runtime adaptation distinct from compile-time staging: one changes a handler after observations, while the other consumes known program syntax before dynamic input arrives.
Both retain the bytecode state contract that makes comparisons with direct interpretation possible.

\section{Exercises}\label{sec:ch31-exercises}

\begin{enumerate}
\item Encode \code{(10 - 4) + (7 - 2)} in the stack bytecode above, compute the depth after each instruction, and give the final stack.\label{ex:ch31-encode-trace}
\item Run Algorithm~\ref{alg:ch31-verify} on a diamond whose two predecessors supply depths $2$ and $3$ to the same join, name the instruction where it reports an error, and propose a one-instruction fix.\label{ex:ch31-depth-join}
\item Add a conditional jump to the machine, state its stack effect and successor rule, and write two tests: one for a jump past the end and one for a depth mismatch after the join.\label{ex:ch31-branch-op}
\end{enumerate}

\section{Further Reading}\label{sec:ch31-reading}

Appel explains virtual machines as translation targets with explicit representations~\cite{appel-1998-modern}.
Morgan relates bytecode layout and calling structure to later optimization choices~\cite{morgan-1998-optimizing}.
Safonov discusses execution trustworthiness when code crosses component boundaries~\cite{safonov-2010-trustworthy}.
The C standard draft defines the sequencing and value rules that any C-facing interface must respect~\cite{wg14-2011-c11}.
